\section{动能定理}

动能定理将合外力做功与质点动能变化联系起来。它是牛顿第二定律的能量形式，在处理已知位移、求速率的问题时尤其方便。

\subsection{质点动能}

质点动能定义为
\[
    \boxed{
    E_k=\frac{1}{2}mv^2
    } .
\]
动能是标量，只与速率大小有关，与速度方向无关。由于速度依赖参考系，动能也依赖参考系。

\subsection{动能定理的推导}

设质点质量为 $m$，在合外力 $\boldsymbol F_{\mathrm{net}}$ 作用下运动。由牛顿第二定律
\[
    \boldsymbol F_{\mathrm{net}}=m\boldsymbol a
\]
可得合外力元功
\[
    \mathrm dW_{\mathrm{net}}
    =
    \boldsymbol F_{\mathrm{net}}\cdot\mathrm d\boldsymbol r
    =
    m\boldsymbol a\cdot\mathrm d\boldsymbol r .
\]
又因为
\[
    \boldsymbol a=\frac{\mathrm d\boldsymbol v}{\mathrm dt},
    \qquad
    \mathrm d\boldsymbol r=\boldsymbol v\,\mathrm dt,
\]
所以
\[
    \mathrm dW_{\mathrm{net}}
    =
    m\frac{\mathrm d\boldsymbol v}{\mathrm dt}\cdot\boldsymbol v\,\mathrm dt
    =
    m\boldsymbol v\cdot\mathrm d\boldsymbol v .
\]
由
\[
    \mathrm d(v^2)
    =
    \mathrm d(\boldsymbol v\cdot\boldsymbol v)
    =
    2\boldsymbol v\cdot\mathrm d\boldsymbol v
\]
可得
\[
    m\boldsymbol v\cdot\mathrm d\boldsymbol v
    =
    \mathrm d\left(\frac{1}{2}mv^2\right).
\]
因此
\[
    \mathrm dW_{\mathrm{net}}=\mathrm dE_k.
\]
从初态到末态积分，得到
\[
    \boxed{
    W_{\mathrm{net}}
    =
    \Delta E_k
    =
    \frac{1}{2}mv_2^2-\frac{1}{2}mv_1^2
    } .
\]

\subsection{使用步骤}

应用动能定理时，可按以下顺序处理：

\begin{enumerate}
    \item 选定研究对象；
    \item 明确初态和末态；
    \item 判断哪些力做功，哪些力不做功；
    \item 写出合外力做功 $W_{\mathrm{net}}$；
    \item 写出动能变化 $\Delta E_k$；
    \item 建立方程
    \[
        W_{\mathrm{net}}=\Delta E_k
    \]
    并求解未知量。
\end{enumerate}

动能定理不直接涉及时间，因此当题目只关心经过某段位移后的速度大小时，通常比逐步求加速度和时间更简洁。

\subsection{与动量定理的区别}

动量定理为
\[
    \boldsymbol I=\Delta\boldsymbol p,
\qquad
    \boldsymbol I=\int\boldsymbol F\,\mathrm dt,
\]
它体现力对时间的累积。动能定理为
\[
    W_{\mathrm{net}}=\Delta E_k,
\qquad
    W=\int\boldsymbol F\cdot\mathrm d\boldsymbol r,
\]
它体现力对空间位移的累积。

因此：

\begin{enumerate}
    \item 冲量是矢量，功是标量；
    \item 动量变化由冲量决定，动能变化由合外力做功决定；
    \item 动量定理保留速度方向信息，动能定理主要处理速率大小变化。
\end{enumerate}

\subsection{例题：粗糙水平面上的制动距离}

质量为 $m$ 的物块以初速度 $v_0$ 在粗糙水平面上滑行，动摩擦因数为 $\mu$。求物块停止前滑行距离 $s$。

\textbf{解析：} 物块受到重力、支持力和滑动摩擦力。重力和支持力都垂直于位移，不做功；滑动摩擦力方向与位移方向相反，大小为
\[
    f_k=\mu mg.
\]
摩擦力做功为
\[
    W_f=-f_ks=-\mu mgs.
\]
物块末速度为零，所以
\[
    \Delta E_k=0-\frac{1}{2}mv_0^2.
\]
由动能定理
\[
    -\mu mgs=-\frac{1}{2}mv_0^2.
\]
解得
\[
    \boxed{
    s=\frac{v_0^2}{2\mu g}
    } .
\]

\subsection{例题：斜向拉力做功}

质量为 $m$ 的物块在水平面上运动，受到大小为 $F$、与水平面夹角为 $\theta$ 的拉力作用，水平位移为 $s$。若忽略摩擦，求拉力做功与物块动能变化。

\textbf{解析：} 位移方向水平，拉力沿位移方向的分量为
\[
    F_x=F\cos\theta.
\]
拉力做功为
\[
    W=Fs\cos\theta.
\]
重力和支持力都垂直于水平位移，不做功。由动能定理
\[
    \Delta E_k=W=Fs\cos\theta.
\]
因此
\[
    \boxed{
    \Delta E_k=Fs\cos\theta
    } .
\]
这说明并非整个拉力都改变动能，只有沿位移方向的分量做功。

\subsection{常用公式}

\[
    \boxed{
    E_k=\frac{1}{2}mv^2
    }
\]
\[
    \boxed{
    W_{\mathrm{net}}=\Delta E_k
    =
    \frac{1}{2}mv_2^2-\frac{1}{2}mv_1^2
    }
\]
\[
    \boldsymbol I=\int\boldsymbol F\,\mathrm dt,\qquad
    W=\int\boldsymbol F\cdot\mathrm d\boldsymbol r.
\]
